\documentclass[a4paper,twoside]{article}

\usepackage{morewrites}

\usepackage[svgnames,dvipsnames,x11names]{xcolor}
\usepackage{graphicx}
\usepackage[english]{babel}
\usepackage[colorlinks=true, linkcolor=NavyBlue, urlcolor=NavyBlue, citecolor=NavyBlue]{hyperref}
\usepackage[a4paper]{geometry}
\usepackage{ts-typesetting}
\usepackage{ts-fonts}
\usepackage{booktabs}
\usepackage{csquotes}
\usepackage{float}
\usepackage{tabularx}
\usepackage[normalem]{ulem}
\usepackage{amsmath}
\usepackage{seqsplit}
\newcolumntype{Y}{>{\centering\arraybackslash}X}
\newcolumntype{Z}{>{\raggedleft\arraybackslash}X}
\usepackage{ts-keystrokes}
\usepackage{ts-callouts}
\usepackage{ts-code}
\usepackage{ts-glossary}
\usepackage{ts-todolist}

\usepackage{lastpage}
\usepackage{refcount}
\makeatletter
\def\ps@plain{
\let\@oddhead\@empty
\let\@evenhead\@empty
\def\@oddfoot{
\hfil
\ifnum\getpagerefnumber{LastPage}>1\relax
\thepage
\fi
\hfil}
\let\@evenfoot\@oddfoot
}
\makeatother

\title{TMark Syntax}
\author{Yves Chevallier}\date{}
\begin{document}
\maketitle
\begin{abstract}
TMark is the Markdown dialect TeXSmith reads: CommonMark, the extensions authors already know from GitHub and MkDocs, and a small closed set of constructs for print --- roles such as \tscodeinline{\{index\}[term]}, containers such as \tscodeinline{::: note}, data fences such as \tscodeinline{yaml table}. This document walks through the syntax construct by construct, each one shown as source and as the typeset result. It doubles as a test sheet for templates.
\end{abstract}
\tableofcontents
\section{Structure}
\subsection{Headings}
Six heading levels, \tscodeinline{\#} to \tscodeinline{\#\#\#\#\#\#}. \LaTeX{} has fewer sectioning commands (\tscodeinline{\textbackslash{}section}, \tscodeinline{\textbackslash{}subsection}, \tscodeinline{\textbackslash{}subsubsection}, \tscodeinline{\textbackslash{}paragraph}, \tscodeinline{\textbackslash{}subparagraph}), so the deepest levels become run-in bold headings.
\begin{tscode}[lang=md, engine=pygments]
\begin{Verbatim}[commandchars=\\\{\},breaklines, breakanywhere, commandchars=\\\{\}]
\PY{g+gh}{\PYZsh{} Heading 1}

Lorem ipsum dolor sit amet.

\PY{g+gu}{\PYZsh{}\PYZsh{} Heading 2}

Lorem ipsum dolor sit amet.

\PY{g+gu}{\PYZsh{}\PYZsh{}\PYZsh{} Heading 3}

Lorem ipsum dolor sit amet.

\PY{g+gu}{\PYZsh{}\PYZsh{}\PYZsh{}\PYZsh{} Heading 4}

Lorem ipsum dolor sit amet.

\PY{g+gu}{\PYZsh{}\PYZsh{}\PYZsh{}\PYZsh{}\PYZsh{} Heading 5}

Lorem ipsum dolor sit amet.
\end{Verbatim}
\end{tscode}
\subsection{Bold}
Strong emphasis with \tscodeinline{**}, typeset as \tscodeinline{\textbackslash{}textbf\{…\}}.
\begin{tscode}[lang=md, engine=pygments]
\begin{Verbatim}[commandchars=\\\{\},breaklines, breakanywhere, commandchars=\\\{\}]
\PY{g+gs}{**text**} or \PY{g+gs}{**t**}ex**t**
\end{Verbatim}
\end{tscode}
\begin{displayquote}
\textbf{text} or \textbf{t}ex\textbf{t}
\end{displayquote}
\subsection{Italic}
Emphasis with \tscodeinline{*} or \tscodeinline{\_}, for terminology or voice.
\begin{tscode}[lang=md, engine=pygments]
\begin{Verbatim}[commandchars=\\\{\},breaklines, breakanywhere, commandchars=\\\{\}]
This \PY{g+ge}{\PYZus{}text\PYZus{}} or \PY{g+ge}{*text*} is italic.
\end{Verbatim}
\end{tscode}
\begin{displayquote}
This \emph{text} or \emph{text} is italic.
\end{displayquote}
\subsection{Nested emphasis}
Bold and italic combine in either spelling.
\begin{tscode}[lang=md, engine=pygments]
\begin{Verbatim}[commandchars=\\\{\},breaklines, breakanywhere, commandchars=\\\{\}]
\PY{g+gs}{**\PYZus{}bold and italic\PYZus{}**}

\PY{g+ges}{***bold and italic***}
\end{Verbatim}
\end{tscode}
\begin{displayquote}
\tslead{\emph{bold and italic}}

\emph{\textbf{bold and italic}}
\end{displayquote}
\subsection{Small capitals}
In CommonMark \tscodeinline{**} and \tscodeinline{\_\_} both mean bold. TMark keeps \tscodeinline{**} for bold and reads \tscodeinline{\_\_} as small capitals.
\begin{tscode}[lang=md, engine=pygments]
\begin{Verbatim}[commandchars=\\\{\},breaklines, breakanywhere, commandchars=\\\{\}]
This is \PY{g+gs}{\PYZus{}\PYZus{}small capitals\PYZus{}\PYZus{}}.
\end{Verbatim}
\end{tscode}
\begin{displayquote}
This is \textsc{small capitals}.
\end{displayquote}
\subsection{Strikethrough}
Marks text as deleted or obsolete.
\begin{tscode}[lang=md, engine=pygments]
\begin{Verbatim}[commandchars=\\\{\},breaklines, breakanywhere, commandchars=\\\{\}]
We \PY{g+gd}{\PYZti{}\PYZti{}do not\PYZti{}\PYZti{}} want this.
\end{Verbatim}
\end{tscode}
\begin{displayquote}
We \sout{do not} want this.
\end{displayquote}
\subsection{Underline}
The \tscodeinline{\{underline\}} role underlines text. PyMdownX's \tscodeinline{\^{}\^{}text\^{}\^{}} spells the same
node, but only with the \tscodeinline{inline.insert} feature enabled; without it the carets
are literal.
\begin{tscode}[lang=md, engine=pygments]
\begin{Verbatim}[commandchars=\\\{\},breaklines, breakanywhere, commandchars=\\\{\}]
\PYZob{}underline\PYZcb{}[text]
\end{Verbatim}
\end{tscode}
\begin{displayquote}
\uline{text}
\end{displayquote}
\subsection{Highlight}
\tscodeinline{==text==} highlights a run of text, as a marker would.
\begin{tscode}[lang=md, engine=pygments]
\begin{Verbatim}[commandchars=\\\{\},breaklines, breakanywhere, commandchars=\\\{\}]
==Cellular respiration is a set of metabolic reactions== that take place
in the cells of organisms. Its ==primary== function is to convert
biochemical energy from nutrients into ==adenosine triphosphate (ATP)==,
and then release waste products.
\end{Verbatim}
\end{tscode}
\begin{displayquote}
\tsmark{Cellular respiration is a set of metabolic reactions} that take place in the
cells of organisms. Its \tsmark{primary} function is to convert biochemical energy
from nutrients into \tsmark{adenosine triphosphate (ATP)}, and then release waste
products.
\end{displayquote}
\subsection{Subscript and superscript}
\tscodeinline{\textasciitilde{}…\textasciitilde{}} lowers, \tscodeinline{\^{}…\^{}} raises. Unicode subscripts and superscripts pass through.
\begin{tscode}[lang=md, engine=pygments]
\begin{Verbatim}[commandchars=\\\{\},breaklines, breakanywhere, commandchars=\\\{\}]
H\PYZti{}2\PYZti{}O, E = mc\PYZca{}2\PYZca{}

H₂O, E = mc²
\end{Verbatim}
\end{tscode}
\begin{displayquote}
H\textsubscript{2}O, E = mc\textsuperscript{2}

H\textsubscript{2}O, E = mc\textsuperscript{2}
\end{displayquote}
\subsection{Smart punctuation}
Straight quotes become typographic quotes, three dots an ellipsis, \tscodeinline{-\allowbreak{}-\allowbreak{}} an en dash and \tscodeinline{-\allowbreak{}-\allowbreak{}-\allowbreak{}} an em dash. A \tscodeinline{-\allowbreak{}-\allowbreak{}-\allowbreak{}} alone on its line is a thematic break, not a dash.
\begin{tscode}[lang=md, engine=pygments]
\begin{Verbatim}[commandchars=\\\{\},breaklines, breakanywhere, commandchars=\\\{\}]
\PYZdq{}Smart quotes\PYZdq{}, ellipses..., en\PYZhy{}dash \PYZhy{}\PYZhy{}, em\PYZhy{}dash \PYZhy{}\PYZhy{}\PYZhy{}

Achilles \PYZhy{}\PYZhy{} the swiftest runner \PYZhy{}\PYZhy{} was fast, but not as fast as the tortoise.
\end{Verbatim}
\end{tscode}
\begin{displayquote}
\enquote{Smart quotes}, ellipses..., en-dash --, em-dash ---

Achilles -- the swiftest runner -- was fast, but not as fast as the tortoise.
\end{displayquote}
\subsection{Inline code and code blocks}
\begin{tscode}[lang=md, engine=pygments]
\begin{Verbatim}[commandchars=\\\{\},breaklines, breakanywhere, commandchars=\\\{\}]
This entry is \PY{l+s+sb}{`inline code`}, or can be in a fenced code block:

\PY{l+s+sb}{```}\PY{l+s+sb}{python}
\PY{n+nb}{print}\PY{p}{(}\PY{l+s+s2}{\PYZdq{}}\PY{l+s+s2}{Hello}\PY{l+s+s2}{\PYZdq{}}\PY{p}{)}
\PY{l+s+sb}{```}
\end{Verbatim}
\end{tscode}
\begin{displayquote}
This entry is \tscodeinline{inline code}, or can be in a fenced code block:
\begin{tscode}[lang=python, engine=pygments]
\begin{Verbatim}[commandchars=\\\{\},breaklines, breakanywhere, commandchars=\\\{\}]
\PY{n+nb}{print}\PY{p}{(}\PY{l+s+s2}{\PYZdq{}}\PY{l+s+s2}{Hello}\PY{l+s+s2}{\PYZdq{}}\PY{p}{)}
\end{Verbatim}
\end{tscode}
\end{displayquote}
\subsection{Links}
External links become clickable links in the PDF. A bare URL is linked as it stands.
\begin{tscode}[lang=md, engine=pygments]
\begin{Verbatim}[commandchars=\\\{\},breaklines, breakanywhere, commandchars=\\\{\}]
The German\PYZhy{}born physicist [\PY{n+nt}{Albert Einstein}](\PY{n+na}{https://en.wikipedia.org/wiki/Albert\PYZus{}Einstein}) is famous for
his Nobel prize and the theory of relativity: https://en.wikipedia.org/wiki/Theory\PYZus{}of\PYZus{}relativity
\end{Verbatim}
\end{tscode}
\begin{displayquote}
The German-born physicist \href{https://en.wikipedia.org/wiki/Albert\_Einstein}{Albert Einstein} is famous for
his Nobel prize and the theory of relativity: \url{https://en.wikipedia.org/wiki/Theory\_of\_relativity}
\end{displayquote}
\subsection{Images}
Local or remote images; a remote one is fetched at build time.
\begin{tscode}[lang=md, engine=pygments]
\begin{Verbatim}[commandchars=\\\{\},breaklines, breakanywhere, commandchars=\\\{\}]
![\PY{n+nt}{Random Picture}](\PY{n+na}{https://picsum.photos/500/200})
\end{Verbatim}
\end{tscode}
\begin{displayquote}
\begin{figure}[H]
\centering
\includegraphics[width=\linewidth]{200.jpg}
\caption{Random Picture}
\end{figure}
\end{displayquote}
\subsection{Bulleted lists}
\begin{tscode}[lang=md, engine=pygments]
\begin{Verbatim}[commandchars=\\\{\},breaklines, breakanywhere, commandchars=\\\{\}]
\PY{k}{\PYZhy{}}\PY{+w}{ }First
\PY{k}{\PYZhy{}}\PY{+w}{ }Second
\PY{+w}{    }\PY{k}{\PYZhy{}}\PY{+w}{ }Subitem
\end{Verbatim}
\end{tscode}
\begin{displayquote}
\begin{itemize}
\item First
\item Second
\begin{itemize}
\item Subitem
\end{itemize}
\end{itemize}
\end{displayquote}
\subsection{Numbered lists}
\begin{tscode}[lang=md, engine=pygments]
\begin{Verbatim}[commandchars=\\\{\},breaklines, breakanywhere, commandchars=\\\{\}]
\PY{k}{1.} Numbered first
\PY{k}{2.} Numbered second
\PY{+w}{    }\PY{k}{1.} Sub\PYZhy{}numbered
        a. Sub\PYZhy{}sub\PYZhy{}numbered
        b. Sub\PYZhy{}sub\PYZhy{}numbered
\PY{+w}{    }\PY{k}{2.} Sub\PYZhy{}numbered
\end{Verbatim}
\end{tscode}
\begin{displayquote}
\begin{enumerate}
\item Numbered first
\item Numbered second
\begin{enumerate}
\item Sub-numbered
a. Sub-sub-numbered
b. Sub-sub-numbered
\item Sub-numbered
\end{enumerate}
\end{enumerate}
\end{displayquote}
\subsection{Task lists}
\begin{tscode}[lang=md, engine=pygments]
\begin{Verbatim}[commandchars=\\\{\},breaklines, breakanywhere, commandchars=\\\{\}]
\PY{k}{\PYZhy{} }\PY{k}{[x]} Done
\PY{k}{\PYZhy{} }\PY{k}{[ ]} To do
\end{Verbatim}
\end{tscode}
\begin{displayquote}
\begin{tstasklist}
\item[\tsdone] Done
\item[\tstodo] To do
\end{tstasklist}
\end{displayquote}
\subsection{Definition lists}
A term, then its definition on a \tscodeinline{:} line.
\begin{tscode}[lang=md, engine=pygments]
\begin{Verbatim}[commandchars=\\\{\},breaklines, breakanywhere, commandchars=\\\{\}]
Term
: Definition of the term

Another term
:   Another definition with a longer description that spans
    multiple lines and
    is indented properly.
\end{Verbatim}
\end{tscode}
\begin{displayquote}
\begin{description}
\item[{Term}] Definition of the term
\item[{Another term}] Another definition with a longer description that spans
multiple lines and
is indented properly.
\end{description}
\end{displayquote}
\subsection{Block quotes}
\begin{tscode}[lang=md, engine=pygments]
\begin{Verbatim}[commandchars=\\\{\},breaklines, breakanywhere, commandchars=\\\{\}]
\PY{k}{\PYZgt{} }\PY{g+ge}{Quote}
\end{Verbatim}
\end{tscode}
\begin{displayquote}
Quote
\end{displayquote}
\subsection{Thematic breaks}
A \tscodeinline{-\allowbreak{}-\allowbreak{}-\allowbreak{}} on its own line. At the top level of a paged document it turns the page; inside a container it draws a rule.
\begin{tscode}[lang=md, engine=pygments]
\begin{Verbatim}[commandchars=\\\{\},breaklines, breakanywhere, commandchars=\\\{\}]
\PYZhy{}\PYZhy{}\PYZhy{}
\end{Verbatim}
\end{tscode}
\begin{displayquote}
\tsrule
\end{displayquote}
\section{Blocks}
\subsection{Tables}
Pipe tables for plain 2-D data; a \tscodeinline{yaml table} fence for spans, grouped headers and footers (see the \tscodeinline{tables} example).
\begin{tscode}[lang=md, engine=pygments]
\begin{Verbatim}[commandchars=\\\{\},breaklines, breakanywhere, commandchars=\\\{\}]
| Column 1 | Column 2 |
| \PYZhy{}\PYZhy{}\PYZhy{}\PYZhy{}\PYZhy{}\PYZhy{}\PYZhy{}\PYZhy{} | \PYZhy{}\PYZhy{}\PYZhy{}\PYZhy{}\PYZhy{}\PYZhy{}\PYZhy{}\PYZhy{} |
| Value 1  | Value 2  |
| Value 3  | Value 4  |
\end{Verbatim}
\end{tscode}
\begin{displayquote}
\begin{center}
\begin{tabularx}{\linewidth}{>{\raggedright\arraybackslash}X>{\raggedright\arraybackslash}X}
\toprule
\textbf{Column 1} & \textbf{Column 2} \\
\midrule
Value 1 & Value 2 \\
Value 3 & Value 4 \\
\bottomrule
\end{tabularx}
\end{center}
\end{displayquote}
\subsection{Footnotes}
A footnote is a note placed at the bottom of the page.
\begin{tscode}[lang=md, engine=pygments]
\begin{Verbatim}[commandchars=\\\{\},breaklines, breakanywhere, commandchars=\\\{\}]
Text with note[\PYZca{}1].

[\PY{n+nl}{\PYZca{}1}]: \PY{n+na}{Here is the note.}
\end{Verbatim}
\end{tscode}
\begin{displayquote}
Text with note\footnote{Here is the note.}.
\end{displayquote}
\subsection{Abbreviations}
An abbreviation definition expands its first occurrence and lists the acronym in the backmatter.
\begin{tscode}[lang=md, engine=pygments]
\begin{Verbatim}[commandchars=\\\{\},breaklines, breakanywhere, commandchars=\\\{\}]
The HTML standard.

*[HTML]: HyperText Markup Language
\end{Verbatim}
\end{tscode}
\begin{displayquote}
The \tsacr{HTML} standard.
\end{displayquote}
\subsection{Callouts}
A \tscodeinline{::: kind} container is a callout: note, tip, warning, danger, and the other kinds of the callout palette. Attributes follow the kind. PyMdownX's \tscodeinline{!!! note} and \tscodeinline{??? warning} spell the same containers.
\begin{tscode}[lang=md, engine=pygments]
\begin{Verbatim}[commandchars=\\\{\},breaklines, breakanywhere, commandchars=\\\{\}]
::: note
This is a note.
:::

::: warning \PYZob{}collapsed=true\PYZcb{}
This is a warning.
:::
\end{Verbatim}
\end{tscode}
\tscodeinline{collapsed=true} creates a collapsible block in \tsacr{HTML} output only.
\begin{tscallout}[kind=note]
This is a note.
\end{tscallout}
\begin{tscallout}[kind=warning, collapsed]
This is a warning.
\end{tscallout}
A title of its own:
\begin{tscode}[lang=md, engine=pygments]
\begin{Verbatim}[commandchars=\\\{\},breaklines, breakanywhere, commandchars=\\\{\}]
::: note \PYZob{}title=Title collapsed=false\PYZcb{}
Collapsible content.
:::
\end{Verbatim}
\end{tscode}
\begin{displayquote}
\begin{tscallout}[kind=note, title={Title}]
Collapsible content.
\end{tscallout}
\end{displayquote}
\subsection{Tabs}
A \tscodeinline{tabs} container holds one \tscodeinline{tab} container per pane. Paged output lays the panes one after the other, each under its title.
\begin{tscode}[lang=md, engine=pygments]
\begin{Verbatim}[commandchars=\\\{\},breaklines, breakanywhere, commandchars=\\\{\}]
:::: tabs
::: tab \PYZob{}title=Windows\PYZcb{}
Windows is a Microsoft operating system.
:::
::: tab \PYZob{}title=Linux\PYZcb{}
Linux is an open\PYZhy{}source operating system.
:::
::::
\end{Verbatim}
\end{tscode}
\begin{tsdiv}{tab}[title=Windows]
Windows is a Microsoft operating system.
\end{tsdiv}
\begin{tsdiv}{tab}[title=Linux]
Linux is an open-source operating system.
\end{tsdiv}
\subsection{Diagrams}
A \tscodeinline{mermaid} fence is rendered to a vector image at build time; \tscodeinline{image} on the fence makes it a figure, with the attributes an image takes.
\begin{tscode}[lang=md, engine=pygments]
\begin{Verbatim}[commandchars=\\\{\},breaklines, breakanywhere, commandchars=\\\{\}]
\PY{l+s+sb}{```}\PY{l+s+sb}{mermaid}\PY{+w}{ }image width=\PYZdq{}20\PYZpc{}\PYZdq{}
\PY{l+s}{graph TD;}
\PY{l+s}{  A\PYZhy{}\PYZhy{}\PYZgt{}B;}
\PY{l+s+sb}{```}
\end{Verbatim}
\end{tscode}
\begin{displayquote}
\begin{figure}[H]
\centering
\includegraphics[width=0.2\linewidth]{<HASH>.pdf}
\end{figure}
\end{displayquote}
\subsection{Math}
Inline math between \tscodeinline{\$…\$}, display math between \tscodeinline{\$\$…\$\$}, in \LaTeX{} syntax on both backends.
\begin{tscode}[lang=md, engine=pygments]
\begin{Verbatim}[commandchars=\\\{\},breaklines, breakanywhere, commandchars=\\\{\}]
The differential form of Maxwell\PYZsq{}s equations:

\PYZdl{}\PYZdl{}
\PYZbs{}begin\PYZob{}aligned\PYZcb{}
\PYZbs{}nabla \PYZbs{}cdot \PYZbs{}mathbf\PYZob{}E\PYZcb{} \PYZam{}= \PYZbs{}frac\PYZob{}\PYZbs{}rho\PYZcb{}\PYZob{}\PYZbs{}varepsilon\PYZus{}0\PYZcb{} \PYZbs{}\PYZbs{}
\PYZbs{}nabla \PYZbs{}cdot \PYZbs{}mathbf\PYZob{}B\PYZcb{} \PYZam{}= 0 \PYZbs{}\PYZbs{}
\PYZbs{}nabla \PYZbs{}times \PYZbs{}mathbf\PYZob{}E\PYZcb{} \PYZam{}= \PYZhy{}\PYZbs{}frac\PYZob{}\PYZbs{}partial \PYZbs{}mathbf\PYZob{}B\PYZcb{}\PYZcb{}\PYZob{}\PYZbs{}partial t\PYZcb{} \PYZbs{}\PYZbs{}
\PYZbs{}nabla \PYZbs{}times \PYZbs{}mathbf\PYZob{}B\PYZcb{} \PYZam{}= \PYZbs{}mu\PY{g+ge}{\PYZus{}0 \PYZbs{}mathbf\PYZob{}J\PYZcb{} + \PYZbs{}mu\PYZus{}}0 \PYZbs{}varepsilon\PYZus{}0
\PYZbs{}frac\PYZob{}\PYZbs{}partial \PYZbs{}mathbf\PYZob{}E\PYZcb{}\PYZcb{}\PYZob{}\PYZbs{}partial t\PYZcb{}
\PYZbs{}end\PYZob{}aligned\PYZcb{}
\PYZdl{}\PYZdl{}
\end{Verbatim}
\end{tscode}
\begin{displayquote}
The differential form of Maxwell's equations:

$$
\begin{aligned}
\nabla \cdot \mathbf{E} &= \frac{\rho}{\varepsilon_0} \\
\nabla \cdot \mathbf{B} &= 0 \\
\nabla \times \mathbf{E} &= -\frac{\partial \mathbf{B}}{\partial t} \\
\nabla \times \mathbf{B} &= \mu_0 \mathbf{J} + \mu_0 \varepsilon_0
\frac{\partial \mathbf{E}}{\partial t}
\end{aligned}
$$
\end{displayquote}
\subsection{Raw passthrough}
A \tscodeinline{latex raw} fence is written to the \LaTeX{} output verbatim; the \tscodeinline{\{raw latex\}(…)} role does the same inside a paragraph. Neither reaches the other backends.
\begin{tscode}[lang=md, engine=pygments]
\begin{Verbatim}[commandchars=\\\{\},breaklines, breakanywhere, commandchars=\\\{\}]
\PY{l+s+sb}{```}\PY{l+s+sb}{latex}\PY{+w}{ }raw
\PY{k}{\PYZbs{}clearpage}
\PY{l+s+sb}{```}
\end{Verbatim}
\end{tscode}
\clearpage
Inline variant:
\begin{tscode}[lang=md, engine=pygments]
\begin{Verbatim}[commandchars=\\\{\},breaklines, breakanywhere, commandchars=\\\{\}]
Insert... \PYZob{}raw latex\PYZcb{}(\PYZbs{}clearpage) anywhere in the paragraph.
\end{Verbatim}
\end{tscode}
Insert... \clearpage anywhere in the paragraph.
\subsection{Including a file}
A fence with \tscodeinline{include} takes its content from a file, here a Python program listed with its highlighting.
\begin{tscode}[lang=md, engine=pygments]
\begin{Verbatim}[commandchars=\\\{\},breaklines, breakanywhere, commandchars=\\\{\}]
\PY{l+s+sb}{```}\PY{l+s+sb}{python}\PY{+w}{ }include=\PYZdq{}hanoi.py\PYZdq{}
\PY{l+s+sb}{```}
\end{Verbatim}
\end{tscode}
\begin{tscode}[lang=python, engine=pygments]
\begin{Verbatim}[commandchars=\\\{\},breaklines, breakanywhere, commandchars=\\\{\}]
\PY{k+kn}{from}\PY{+w}{ }\PY{n+nn}{\PYZus{}\PYZus{}future\PYZus{}\PYZus{}}\PY{+w}{ }\PY{k+kn}{import} \PY{n}{annotations}

\PY{k+kn}{from}\PY{+w}{ }\PY{n+nn}{collections}\PY{n+nn}{.}\PY{n+nn}{abc}\PY{+w}{ }\PY{k+kn}{import} \PY{n}{Iterable}

\PY{n}{Move} \PY{o}{=} \PY{n+nb}{tuple}\PY{p}{[}\PY{n+nb}{int}\PY{p}{,} \PY{n+nb}{str}\PY{p}{,} \PY{n+nb}{str}\PY{p}{]}

\PY{k}{def}\PY{+w}{ }\PY{n+nf}{tower\PYZus{}of\PYZus{}hanoi}\PY{p}{(}\PY{n}{n}\PY{p}{:} \PY{n+nb}{int}\PY{p}{,} \PY{n}{source}\PY{p}{:} \PY{n+nb}{str}\PY{p}{,} \PY{n}{destination}\PY{p}{:} \PY{n+nb}{str}\PY{p}{,} \PY{n}{auxiliary}\PY{p}{:} \PY{n+nb}{str}\PY{p}{)} \PY{o}{\PYZhy{}}\PY{o}{\PYZgt{}} \PY{n+nb}{list}\PY{p}{[}\PY{n}{Move}\PY{p}{]}\PY{p}{:}
\PY{+w}{    }\PY{l+s+sd}{\PYZdq{}\PYZdq{}\PYZdq{}Compute the move sequence for the Tower of Hanoi puzzle.\PYZdq{}\PYZdq{}\PYZdq{}}
    \PY{n}{moves}\PY{p}{:} \PY{n+nb}{list}\PY{p}{[}\PY{n}{Move}\PY{p}{]} \PY{o}{=} \PY{p}{[}\PY{p}{]}

    \PY{k}{def}\PY{+w}{ }\PY{n+nf}{\PYZus{}solve}\PY{p}{(}\PY{n}{disks}\PY{p}{:} \PY{n+nb}{int}\PY{p}{,} \PY{n}{start}\PY{p}{:} \PY{n+nb}{str}\PY{p}{,} \PY{n}{end}\PY{p}{:} \PY{n+nb}{str}\PY{p}{,} \PY{n}{spare}\PY{p}{:} \PY{n+nb}{str}\PY{p}{)} \PY{o}{\PYZhy{}}\PY{o}{\PYZgt{}} \PY{k+kc}{None}\PY{p}{:}
        \PY{k}{if} \PY{n}{disks} \PY{o}{==} \PY{l+m+mi}{0}\PY{p}{:}
            \PY{k}{return}
        \PY{n}{\PYZus{}solve}\PY{p}{(}\PY{n}{disks} \PY{o}{\PYZhy{}} \PY{l+m+mi}{1}\PY{p}{,} \PY{n}{start}\PY{p}{,} \PY{n}{spare}\PY{p}{,} \PY{n}{end}\PY{p}{)}
        \PY{n}{moves}\PY{o}{.}\PY{n}{append}\PY{p}{(}\PY{p}{(}\PY{n}{disks}\PY{p}{,} \PY{n}{start}\PY{p}{,} \PY{n}{end}\PY{p}{)}\PY{p}{)}
        \PY{n}{\PYZus{}solve}\PY{p}{(}\PY{n}{disks} \PY{o}{\PYZhy{}} \PY{l+m+mi}{1}\PY{p}{,} \PY{n}{spare}\PY{p}{,} \PY{n}{end}\PY{p}{,} \PY{n}{start}\PY{p}{)}

    \PY{n}{\PYZus{}solve}\PY{p}{(}\PY{n}{n}\PY{p}{,} \PY{n}{source}\PY{p}{,} \PY{n}{destination}\PY{p}{,} \PY{n}{auxiliary}\PY{p}{)}
    \PY{k}{return} \PY{n}{moves}

\PY{k}{def}\PY{+w}{ }\PY{n+nf}{render\PYZus{}solution}\PY{p}{(}\PY{n}{moves}\PY{p}{:} \PY{n}{Iterable}\PY{p}{[}\PY{n}{Move}\PY{p}{]}\PY{p}{)} \PY{o}{\PYZhy{}}\PY{o}{\PYZgt{}} \PY{n+nb}{str}\PY{p}{:}
\PY{+w}{    }\PY{l+s+sd}{\PYZdq{}\PYZdq{}\PYZdq{}Return a human\PYZhy{}readable description of the move sequence.\PYZdq{}\PYZdq{}\PYZdq{}}
    \PY{k}{return} \PY{l+s+s2}{\PYZdq{}}\PY{l+s+se}{\PYZbs{}n}\PY{l+s+s2}{\PYZdq{}}\PY{o}{.}\PY{n}{join}\PY{p}{(}
        \PY{l+s+sa}{f}\PY{l+s+s2}{\PYZdq{}}\PY{l+s+s2}{Move disk }\PY{l+s+si}{\PYZob{}}\PY{n}{disk}\PY{l+s+si}{\PYZcb{}}\PY{l+s+s2}{ from source }\PY{l+s+si}{\PYZob{}}\PY{n}{start}\PY{l+s+si}{\PYZcb{}}\PY{l+s+s2}{ to destination }\PY{l+s+si}{\PYZob{}}\PY{n}{end}\PY{l+s+si}{\PYZcb{}}\PY{l+s+s2}{\PYZdq{}} \PY{k}{for} \PY{n}{disk}\PY{p}{,} \PY{n}{start}\PY{p}{,} \PY{n}{end} \PY{o+ow}{in} \PY{n}{moves}
    \PY{p}{)}
\end{Verbatim}
\end{tscode}
\section{Inline}
\subsection{Emoji}
Shortcodes and Unicode emoji are typeset with the emoji font of the \tscodeinline{fonts.emoji} setting --- OpenMoji black unless told otherwise.
\begin{tscode}[lang=md, engine=pygments]
\begin{Verbatim}[commandchars=\\\{\},breaklines, breakanywhere, commandchars=\\\{\}]
:smile:
:heart:
\end{Verbatim}
\end{tscode}
\begin{displayquote}
\tsemoji{😄}
\tsemoji{❤️}
\end{displayquote}
\begin{tscode}[lang=md, engine=pygments]
\begin{Verbatim}[commandchars=\\\{\},breaklines, breakanywhere, commandchars=\\\{\}]
😊 🚀 🍕 🎉 🐍 🌍 💻
📚 🎨 👽 👋 🤖 🦄 🧠
\end{Verbatim}
\end{tscode}
\begin{displayquote}
\tsemoji{😊} \tsemoji{🚀} \tsemoji{🍕} \tsemoji{🎉} \tsemoji{🐍} \tsemoji{🌍} \tsemoji{💻}
\tsemoji{📚} \tsemoji{🎨} \tsemoji{👽} \tsemoji{👋} \tsemoji{🤖} \tsemoji{🦄} \tsemoji{🧠}
\end{displayquote}
\subsection{Keys}
\tscodeinline{++key+key++} sets keyboard shortcuts as key caps.
\begin{tscode}[lang=md, engine=pygments]
\begin{Verbatim}[commandchars=\\\{\},breaklines, breakanywhere, commandchars=\\\{\}]
++Ctrl+C++
\end{Verbatim}
\end{tscode}
\begin{displayquote}
\tskeys{Ctrl,C}
\end{displayquote}
\subsection{Progress bars}
\begin{tscode}[lang=markdown, engine=pygments]
\begin{Verbatim}[commandchars=\\\{\},breaklines, breakanywhere, commandchars=\\\{\}]
[=25\PYZpc{} \PYZdq{}Research\PYZdq{}]
[=50\PYZpc{} \PYZdq{}Implementation\PYZdq{}]
[=75\PYZpc{} \PYZdq{}Review\PYZdq{}]
[=100\PYZpc{} \PYZdq{}Launch\PYZdq{}]\PYZob{}.thin\PYZcb{}
\end{Verbatim}
\end{tscode}
\tsprogress{0.25}{Research}
\tsprogress{0.5}{Implementation}
\tsprogress{0.75}{Review}
\tsprogress[thin]{1}{Launch}
\subsection{Escaping}
A backslash makes any punctuation literal, so a \tscodeinline{*}, a \tscodeinline{\_} or a \tscodeinline{\{} can be written as it is.
\begin{tscode}[lang=md, engine=pygments]
\begin{Verbatim}[commandchars=\\\{\},breaklines, breakanywhere, commandchars=\\\{\}]
\PYZbs{}*not emphasis\PYZbs{}*, a literal \PYZbs{}\PYZob{}brace\PYZbs{}\PYZcb{} and a \PYZbs{}\PYZsh{} that is no heading.
\end{Verbatim}
\end{tscode}
\begin{displayquote}
*not emphasis*, a literal \{brace\} and a \# that is no heading.
\end{displayquote}
\section{Front matter}
A YAML block at the top of the document carries the metadata and the \tscodeinline{press} settings: the template, its attributes, the slots, the fonts.
\begin{tscode}[lang=md, engine=pygments]
\begin{Verbatim}[commandchars=\\\{\},breaklines, breakanywhere, commandchars=\\\{\}]
\PYZhy{}\PYZhy{}\PYZhy{}
title: My document
author: Alice
press:
  template: article
\PY{g+gu}{  toc: true}
\PY{g+gu}{\PYZhy{}\PYZhy{}\PYZhy{}}

\PY{g+gh}{\PYZsh{} Title}
\end{Verbatim}
\end{tscode}
\printglossary[type=\acronymtype]
\end{document}
